% !TEX root = ../algebra_02.tex

\section{Левые и правые смежные классы. Индекс подгруппы}

Пусть $H < G$. Введем на группе $G$ отношение $\sim$: $g' \sim g \Leftrightarrow g' = gh$ для некоторого $h \in H$. Это отношение является отношением эквивалентности.

\begin{Statement}{$\sim$~--- отношение эквивалентности}{}
	
	Проверим:
	
	\begin{enumerate}
		\item $g \sim g \st g = g\cdot e, \qquad e \in H$ (т.к. $H < G$)
		\item $g' \sim g \Rightarrow g \sim g' \st g' = gh \Rightarrow g = g'h^{-1}$ (т.к. $h\in H \Rightarrow h^{-1} \in H$)
		\item $g_1 \sim g_2, \ \ g_2 \sim g_3 \Rightarrow g_1 \sim g_3 \st$
		\[ g_1 = g_2h, \ \ g_2 = g_3h', \qquad h, h' \in H\]
		\[ g_1 = g_3\cdot h\cdot h', \qquad h\cdot h' = \tilde{h} \in H\]
		\[ g_1 = g_3 \tilde{h} \Leftrightarrow g_1 \sim g_3 \]
	\end{enumerate}
\end{Statement}

\begin{Definition}{Левый смежный класс}{}
	Класс эквивалентности элемента $g$ по этому отношению:
	\[
	gH := \{ gh \mid h \in H \}
	\]
	называется левым смежным классом элемента $g$ по подгруппе $H$.
\end{Definition}

Множество всех левых смежных классов обозначается $G/H := \{ gH \mid g \in G \}$.

\begin{Definition}{Индекс подгруппы}{}
	Количество левых классов смежности в группе называется индексом подгруппы $H$ в $G$ и обозначается $(G:H)$ или $|G/H|$.
\end{Definition}

\begin{Definition}{Правый смежный класс}{}
	Аналогично, правый класс смежности элемента $g$ по $H$ определяется как $Hg$. Множество правых смежных классов обозначается $H \backslash G$.
\end{Definition}

Количество левых и правых смежных классов всегда совпадает:
\begin{Proposition}{}{}
	$|H \backslash G| = (G:H) = |G \slash H|$.
\end{Proposition}
\begin{proof}
	Рассмотрим отображение $\alpha\colon G \slash H \to H \backslash G$, заданное правилом $gH \mapsto Hg^{-1}$. Оно корректно, так как класс $p^{-1} = H^{-1}g^{-1} = Hg^{-1}$ действительно является правым смежным классом.
	Существует обратное отображение $\beta \colon Hg \mapsto g^{-1}H$, откуда следует, что $\alpha$ --- биекция, а значит множества равномощны.
\end{proof}
